\section{指令微调（Supervised Fine-Tuning）}

\subsection{从预训练到指令遵循}

指令微调（Instruction Fine-tuning，也称 Supervised Fine-tuning / SFT）的核心思想非常直接：收集大量的\textbf{（指令，期望输出）对}，然后用标准的 next-token prediction 损失来微调预训练模型。

\begin{figure}[H]
\centering
\includegraphics[width=0.95\textwidth]{slides-images/slide-030.jpg}
\caption{指令微调的范式：收集多种任务的指令-输出对，在预训练模型上进行微调\protect\footnotemark}
\end{figure}
\footnotetext{Slides 第31页。}

\begin{importantbox}{指令微调的核心流程}
\begin{enumerate}
    \item \textbf{收集数据}：从多种任务（QA、翻译、摘要、代码、推理等）收集指令和期望输出的配对
    \item \textbf{微调模型}：用这些配对数据对预训练模型进行 supervised fine-tuning
    \item \textbf{评估}：在未见过的任务上测试模型的泛化能力
\end{enumerate}
\end{importantbox}

这里的关键创新不是技术本身——它只是标准的微调。关键在于\textbf{数据的规模和多样性}：不再是针对单一任务微调，而是同时在成千上万种不同任务上微调。

\begin{figure}[H]
\centering
\includegraphics[width=0.95\textwidth]{slides-images/slide-031.jpg}
\caption{指令微调数据的规模演进：从单一任务到数千任务、数百万样本\protect\footnotemark}
\end{figure}
\footnotetext{Slides 第32页。}

\subsection{规模效应：更大的模型从指令微调中获益更多}

一个令人鼓舞的趋势是：模型越大，从指令微调中获得的提升就越大。

\begin{figure}[H]
\centering
\includegraphics[width=0.95\textwidth]{slides-images/slide-034.jpg}
\caption{FLAN-T5 实验：模型越大，指令微调带来的增益越大（从 +6.1 到 +26.6）\protect\footnotemark}
\end{figure}
\footnotetext{Slides 第35页。}

以 FLAN-T5 系列为例：
\begin{itemize}
    \item T5-Small（60M 参数）：指令微调提升 +6.1 分
    \item T5-XXL（11B 参数）：指令微调提升 \textbf{+26.6} 分
\end{itemize}

这意味着预训练的``基础能力''和指令微调的``对齐能力''之间存在正向的协同效应。

\begin{figure}[H]
\centering
\includegraphics[width=0.95\textwidth]{slides-images/slide-035.jpg}
\caption{指令微调前后的对比：微调前模型无法正确遵循``Let's think step by step''指令，微调后能够清晰地推理并给出答案\protect\footnotemark}
\end{figure}
\footnotetext{Slides 第36页。}

\subsection{评估：MMLU 基准}

\begin{figure}[H]
\centering
\includegraphics[width=0.95\textwidth]{slides-images/slide-033.jpg}
\caption{MMLU 基准测试：涵盖天文学、生物学等 57 个学科的多选题\protect\footnotemark}
\end{figure}
\footnotetext{Slides 第34页。}

MMLU（Massive Multitask Language Understanding）是一个包含 57 个学科的多选题基准。90\% 被视为``人类水平''的标准线。从 GPT-2 时代到 Gemini，模型在此基准上取得了巨大进步。

\begin{warningbox}{评估的陷阱}
大规模评估面临一个根本性问题：\textbf{测试集泄露}。当预训练数据包含数万亿 token 时，很难保证测试题目没有出现在训练数据中。此外，当模型错误时，往往是因为题目本身含糊不清。这使得基准分数的含义变得模糊——90\% 以上的得分未必代表真正的``人类水平理解''。
\end{warningbox}

\subsection{指令微调的实用经验}

近期研究揭示了几个重要的实用发现：

\begin{figure}[H]
\centering
\includegraphics[width=0.95\textwidth]{slides-images/slide-037.jpg}
\caption{指令微调数据集的生态系统：大量开源数据集涌现\protect\footnotemark}
\end{figure}
\footnotetext{Slides 第38页。}

\textbf{1. 模型蒸馏}：可以使用强大的模型（如 GPT-4）生成指令微调数据，用于训练较小的开源模型。这大大降低了数据收集成本。

\textbf{2. 数据质量 vs 数量}：``LIMA: Less Is More for Alignment'' 论文表明，仅用约 1000 条高质量样本就能达到不错的指令遵循能力。数据质量可能比数量更重要。

\textbf{3. 众包数据}：Open Assistant 等项目展示了社区驱动的数据收集方式的可行性。

\begin{knowledgebox}{代码数据的特殊价值}
课堂讨论中提到一个有趣的观点：代码数据天然具有``逐步推理''的结构——每一行代码都是将问题分解为更小步骤的过程。因此，在预训练混合数据中增加代码数据的比例，已被证明能提升模型的推理能力，即使在非代码任务上也是如此。
\end{knowledgebox}

\subsection{指令微调的三大局限}

指令微调虽然有效，但存在三个根本性局限，正是这些局限推动了 RLHF 和 DPO 的发展：

\begin{figure}[H]
\centering
\includegraphics[width=0.95\textwidth]{slides-images/slide-041.jpg}
\caption{指令微调的三大局限\protect\footnotemark}
\end{figure}
\footnotetext{Slides 第42页。}

\textbf{局限 1：数据收集昂贵。} 预训练可以直接爬取互联网数据，但指令微调需要人工编写高质量的回答。随着任务复杂度提升（如物理学博士级别的问题），这变得越来越昂贵。

\textbf{局限 2：开放性任务无标准答案。} 对于创意写作、故事生成等任务，不存在唯一的``正确答案''。这使得收集训练数据变得困难——你该让标注者写什么样的回答？

\textbf{局限 3：Token 级损失无法区分错误严重程度。} Next-token prediction 损失函数对所有 token 级错误一视同仁。但将 ``Avatar'' 描述为 ``adventure show''（可以接受）和描述为 ``musical''（严重错误）的代价是完全不同的。

\begin{figure}[H]
\centering
\includegraphics[width=0.95\textwidth]{slides-images/slide-042.jpg}
\caption{并非所有错误都一样严重：``adventure'' vs ``musical'' 的错误程度截然不同，但 token 级损失无法区分\protect\footnotemark}
\end{figure}
\footnotetext{Slides 第43页。}

\begin{warningbox}{更深层的问题：目标函数不匹配}
指令微调最根本的问题是：我们的\textbf{真正目标}是让模型生成人类喜欢的输出，但\textbf{优化目标}仍然是 token 级的对数似然。这两者之间存在本质的不匹配。此外，随着模型变得越来越强大，人类标注者的回答质量可能反而成为瓶颈——模型生成的回答可能已经比人类标注者的更好。
\end{warningbox}

\subsection{本章小结}

\begin{itemize}
    \item 指令微调通过在大量（指令，回答）对上微调预训练模型，使其学会遵循指令
    \item 更大的模型从指令微调中获益更多，体现了预训练和对齐之间的正向协同
    \item 数据质量可能比数量更重要；可以用强大模型生成训练数据
    \item 三大局限——高昂的数据收集成本、开放任务缺乏标准答案、损失函数无法区分错误严重程度——推动了基于偏好的优化方法的发展
\end{itemize}

%% ========== Part 4: RLHF ==========

